\section{Notations, Conventions}
\uuline{数について}: $\N$ は正の整数全体の集合、$\Z$ は整数全体の集合、$\Q$ は有理数全体の集合、$\R$ は実数全体の集合、$\C$ は複素数全体の集合を表す。また、$\Nzero$ は非負整数全体の集合を表す。

$\Zhat$ にて、$\Z$ の副有限完備化をあらわすものとする。

$\Primes$ にて、素数全体の集合をあらわすものとする。

$p \in \Primes$ について、$\Zp$ は $\Z$ の $p$-進完備化をあらわすものとする。また、$\Qp$ は $\Zp$ の商体をあらわすものとする。また、$p \in \Primes$ について、$\Fp$ は体 $\Z / p\Z$ をあらわすものとする。また、文脈上 $\ell$ とは、なんらかの体 (環) を係数とする状況において、その標数あるいは剰余標数、あるいは特別な素数が $p$ であるとき、その $p$ と異なる素数をあらわすものとする。

\uuline{体の種類について}: $\Q$ の有限次拡大について、これを数体 (number field) とよび、NF と省略する。なんらかの素数 $p$ について $\Fp$ の有限次拡大とあらわされる体について、これを有限体 (finite field) とよび、FF と省略する。なんらかの素数 $p$ について $\Fp$ の代数拡大とあらわされる体について、これを quasi-finite field とよび、QFF と省略する。なんらかの素数 $p$ について $\Qp$ の有限次拡大とあらわされる体について、これを混標数局所体 (mixed-characteristic local field) あるいはこの $p$ について $p$-進局所体 ($p$-adic local field) とよび、MLF あるいは MLF$_p$ と省略する。なんらかの素数 $p$ について $\Qp$ の代数拡大とあらわされる体について、これを mixed-characteristic quasi-local field あるいはこの $p$ について $p$-adic quasi-local field とよび、MQLF あるいは MQLF$_p$ と省略する。また、なんらかの素数 $p$ について $\Fp$ 上の $1$ 変数関数体の代数拡大と同型であるような体について、これを正標数大域体 (positive-characteristic global field) とよび、PGF と省略する。なんらかの素数 $p$ について $\Fp$ 上の $1$ 変数ローラン級数のなす体の有限次拡大と同型であるような体について、これを正標数局所体 (positive-characteristic local field) とよび、PLF と省略する。

\uuline{位相群について}: $G$ を Hausdorff 位相群とする。また $H \subset G$ を閉部分群とする。このとき、$Z_G(H)$ を $H$ の $G$ における中心化群 (centralizer) とする。また、$N_G(H)$ を $H$ の $G$ における正規化群 (normalizer) とする。また、$C_G(H)$ を $H$ の $G$ における commensurator とする。$N_G(H) = H$ (resp., $C_G(H) = H$) であるとき、$H$ は $G$ において normally terminal (resp., commensurably terminal) であるという。ここで、$C_G(H)$ は必ずしも閉部分群であるとは限らない。

$G^\ab$ を $G$ のアーベル化 (abelianization) とする。また、$G^\abtor$ を $G^\ab$ をその torsion 部分群の閉包で割ったものとする。$G$ の副有限完備化を $\profcompl{G}$ と表記する。また、$G^p$ を $G$ の pro-$p$ 完備化、$G^\solv$ を $G$ の pro-solvable 完備化とする。$\Aut(G)$ を $G$ の自己同型群とする。また、$\Out(G)$ を $G$ の外部自己同型群とする。$\Inn(G)$ を $G$ の内部自己同型群とする。$G$ が位相的有限生成であるとき、$G$ が特性部分群による開基をもつことを鑑みれば、$\Aut(G)$ に自然な副有限群の構造を与えることができる。この状況においては、$\Aut(G)$ はここに定めた位相を備えるものとする。$\Out(G)$ についても同様の位相を備えるものとする。

$G$ が center-free であるとは、$G$ の中心 $\Z(G)$ が自明であることをいう。$G$ が center-free であるとき、$\Inn(G)$ は $G$ と同型である --- したがって、次の完全列が存在する。
\[
1 \to G \simeq \Inn(G) \to \Aut(G) \to \Out(G) \to 1.
\]
このとき、群準同型 $Q \to \Out(G)$ に対して、この準同型による引き戻しとして、$1 \to G \to E \to Q \to 1$ なる完全列を構成することができる。

\uuline{代数的スタック・log scheme について}: 代数的スタックの coarse space に関する議論は、\cite[Chapter I, \S 4.10]{FC} を参照されたい。代数的スタックが scheme-like であるとは、スキームであることをいう。また、代数的スタックが generically scheme-like であるとは、open dense な部分スタックがスキームであることをいう。

log scheme に関する事項は、\cite{Kato} を基本的な文献とする。log scheme の内部 (interior) とは、その log 構造が自明であるような最大の開部分スキームをさしていう。

\uuline{曲線について}:
